% TT30041C
% Verletzungen mit chemischen Substanzen
% 14.0
% 2013-11-30
:ref:`m-chemische-verletzungen`
\begin{enumerate}
  -   Auf Selbsschutz achten!
  \begin{itemize}
    -   Patienten ggfs. retten, aber nur wenn es der
    Eigenschutz zulässt.
  \end{itemize}
  -   Gefahr erkennen: \EE{Welcher Stoff?}  Expertenrat einholen!
  -   Wenn erforderlich: Vorgehen analog zu Abschnitt
  \enquote{Gefahrengutunfall}, GAS-Regel
  (\prettyref{topic:gefahrengutunfall})
  -   Weitere Opfer ausschließen (lassen)
  -   Sicherung der Vitalfunktionen, Beurteilung ob der
  Patient vital bedroht ist.
  -   |TxBeiVit| |TxMassVitMKBes|
  \begin{itemize}
    -   Lagerung: situationsgerecht.
  \end{itemize}
  -   Entfernung kontaminierter Kleidung
  -   Spülung der betroffenen Hautpartien mit Wasser  (richtig abrinnen
  lassen)
  -   Auge: Spülung mit Wasser / Ringerlösung (nach außen abrinnen
  lassen!),  Kontaktlinsen entfernen
  -   Transportentscheidung: Ab einer Ausdehnung von
  15% (Erwachsener) bzw. 10% (Kind) Spezialabteilung
  für Verbrennungsverletzungen. Bei Augenverletzungen Abt.
  für Augenheilkunde.
\end{enumerate}